% QMB_n08_qubit_formalismus.tex — Chapter 8: Qubit Formalism
% Substantive basis: Doc. 175 (March 2026), Doc. 034, Doc. 230; own presentation
\chapter{More Than Two States: The Qubit Geometrically}
\label{ch:qubit}

\section{What Lies Between $|0\rangle$ and $|1\rangle$?}

Spin-up and spin-down are two points. Two points are uniquely
labeled --- one bit of information. But a qubit is not a bit. The
actual power of the qubit lies not in the fact that it can be $|0\rangle$
\emph{or} $|1\rangle$; a classical coin can do that too.
It lies in the fact that it can superpose both in a continuous ratio:
\begin{equation}
  |\psi\rangle = \cos\tfrac{\theta}{2}\,|0\rangle
  + e^{i\phi}\sin\tfrac{\theta}{2}\,|1\rangle.
\end{equation}
Two parameters, the mixing angle $\theta\in[0,\pi]$ and the relative
phase $\phi\in[0,2\pi)$, parameterize a spherical surface: the
Bloch sphere. On it there are infinitely many points --- not only the
poles $|0\rangle$ and $|1\rangle$, but the entire equator, all
latitude circles, every point in between. Each is a valid,
physically distinguishable state.

\begin{keypoint}[A qubit encodes not 2 states, but infinitely many]
The two classical poles are special cases of a continuous
family. What distinguishes a qubit from a bit is not the number of
discrete states, but the continuous geometry of the
state space. Readout couples to the configuration and drives it
into the nearest stable extremal state. The result is, in the FFGFT view, deterministic --- it depends on the exact initial state,
which we do not fully know before the measurement. The states
themselves are discrete; the continuity lies in the phase space \emph{before}
the measurement.
\end{keypoint}

\section{The Cylindrical Phase Space}

The Bloch sphere is the standard representation. The FFGFT has a
physically more direct alternative: the cylindrical phase space with
three coordinates $(z,r,\theta)$. Here $z\in[-1,+1]$ is the height
along the cylinder axis --- the proportion of the torsion configuration in
the direction of $|0\rangle$ versus $|1\rangle$; $r\in[0,1]$ the radius in the
cross-section --- the strength of the intermediate torsion, i.e., how far the
state is from both poles; $\theta\in[0,2\pi)$ the azimuthal angle
--- the relative phase between the basis components. The
relationship to the Bloch sphere is direct: $z = \cos\theta_{\mathrm{Bloch}}$,
$r = |\sin\theta_{\mathrm{Bloch}}|$; the formal bridge is in
Chapter~\ref{ch:hilbert}.

The cylinder is not a geometric aesthetic device: It makes
gate operations into geometric transformations in real space.

\begin{center}
\small\renewcommand{\arraystretch}{1.4}
\begin{tabular}{lp{2.4cm}p{4.2cm}}
\toprule
\textbf{Gate} & \textbf{Matrix} & \textbf{Cylinder Transformation} \\
\midrule
Hadamard $H$ & $\frac{1}{\sqrt2}\begin{pmatrix}1&1\\1&-1\end{pmatrix}$ & Swap $z\leftrightarrow r$, rotate phase by $\pi/2$: moves pole to equator \\
Phase $Z$ & $\begin{pmatrix}1&0\\0&-1\end{pmatrix}$ & Add $\pi$ to phase $\theta$: rotation about the cylinder axis \\
Bit-Flip $X$ & $\begin{pmatrix}0&1\\1&0\end{pmatrix}$ & $(z,r)$ rotation by $\alpha = \pi K_{\mathrm{frak}}$ \\
\bottomrule
\end{tabular}
\end{center}

In the cylindrical formalism, every gate operation is a geometric
transformation of the torsion vector. The Hadamard gate does not rotate a
matrix --- it tilts the configuration from the longitudinal axis onto the
transverse plane. The $Z$ gate rotates the winding about its own
axis. The factor $K_{\mathrm{frak}}$ in the $X$ gate is not a subsequently
inserted correction factor; it follows from the fractal
space-time geometry and means: Every qubit rotation is slightly smaller
than $\pi$, because space-time itself deviates minimally from flat geometry. This is a testable prediction: All $\pi$ gates should
systematically deviate from the ideal value by $\Delta\alpha\approx0{,}013\,\%$.\status{S}

\section{Quantum Registers: Exponential Growth Through Combination}

A single qubit has infinitely many intermediate states, but on
readout always yields one of the two poles. The actual power comes
through combination.

\textbf{Product states: additive classical.} Two independent qubits
together have four basis states. The general state without
entanglement is $(\alpha_A|0\rangle+\beta_A|1\rangle)\otimes
(\alpha_B|0\rangle+\beta_B|1\rangle)$ --- two separate Bloch spheres,
no gain over two independent qubits. The state space is additive:
$2+2 = 4$ parameters.

\textbf{Entangled states: multiplicative growth.} As soon as the
qubits are entangled, the joint state can no longer be written as a
product. The general two-qubit state
\begin{equation}
  |\Psi\rangle = \alpha|00\rangle+\beta|01\rangle+\gamma|10\rangle+\delta|11\rangle
\end{equation}
requires four complex amplitudes, i.e., $2\cdot4-2 = 6$ real parameters.
Entangled qubits carry joint information that is in none of the
individual qubits. For $n$ qubits the Hilbert space grows exponentially,
$d(n) = 2^n$, with $2\cdot2^n-2$ real parameters:

\begin{center}
\small
\begin{tabular}{rrrl}
\toprule
$n$ Qubits & Basis States & Real Parameters & Classical \\
\midrule
1 & 2 & 2 & 1 Bit \\
2 & 4 & 6 & 2 Bits \\
3 & 8 & 14 & 3 Bits \\
10 & 1024 & 2046 & 10 Bits \\
50 & $10^{15}$ & $2\times10^{15}$ & 50 Bits \\
300 & $10^{90}$ & $2\times10^{90}$ & $\approx$ atoms in the universe \\
\bottomrule
\end{tabular}
\end{center}

That is the actual quantum advantage: Not a single qubit is
powerful, but the collective growth of an entangled register.

\textbf{Bell states.} The four maximally entangled
two-qubit states form an orthonormal basis of the entangled
subspace:
\begin{align}
  |\Phi^\pm\rangle &= \tfrac{1}{\sqrt2}(|00\rangle\pm|11\rangle),\\
  |\Psi^\pm\rangle &= \tfrac{1}{\sqrt2}(|01\rangle\pm|10\rangle).
\end{align}
If one measures qubit~A, the result of qubit~B is determined, no matter how
far away. In the FFGFT picture, a Bell state is not action at a distance
between two qubits, but a topologically connected
torsion configuration across two resonators: Both share a
common torsion knot, whose resolution upon any local measurement determines
both ends simultaneously. This is not a signal with
superluminal speed, but a common geometric state
that was established during preparation (Chapter~\ref{ch:verschraenkung}).
The earlier damping formula for the Bell correlation, which lowered the CHSH value
by $\sim10^{-4}$, is replaced by the two-observable reading from
Chapter~\ref{ch:chsh_exp}.

\textbf{Hardware validation.} The FFGFT predictions for Bell states
were tested on IBM Brisbane and IBM Sherbrooke (127 qubits each, Heavy-Hex
lattice): Bell state fidelity $97{,}17\,\%$; FFGFT algorithm
and standard QM agree within measurement accuracy;
repeat variance $10^{-4}$ to $10^{-3}$; deviations in $P(00)$ and
$P(11)$ of $1$--$5\,\%$ are hardware noise.\status{K} The
$\xi$ deviation from the ideal Bell correlation lies below the
noise threshold and could not be isolated. The complete
Python code is available in the repository.

\section{Three Ways to More States}

From resonator physics, three fundamentally different
strategies emerge to encode more than two states:

\begin{center}
\small\renewcommand{\arraystretch}{1.3}
\begin{tabular}{p{1.6cm}p{2.3cm}p{1.6cm}p{2.5cm}}
\toprule
\textbf{Strategy} & \textbf{More States Through} & \textbf{Growth} & \textbf{FFGFT Meaning} \\
\midrule
Superposition (1 qubit) & continuous superposition & $\infty$, but 1 bit on measurement & intermediate torsion \\
Register ($n$ qubits) & entanglement & $2^n$ & topologically connected knots \\
Qudit ($d>2$) & higher levels, orbitals & $d$ per object & $d$ stable configurations in one resonator \\
\bottomrule
\end{tabular}
\end{center}

No strategy dominates everywhere. Register entanglement scales most
strongly, but requires $n$ physical objects and error-corrected
entanglement operations. Qudit encoding gains information per object
($\log_2d>1$); two ququarts ($d = 4$) carry $4^2 = 16$ states, four
qubits $2^4 = 16$ --- mathematically equivalent, physically differently
robust. Superposition is always there: It is the basis of
interference, the actual driving force of quantum algorithms.

\section{Quantum Interference: Why Superposition Is Useful}

Superposition alone would be a mere superposition; readout would always
yield one of the two poles. What makes superposition a resource is
interference: Amplitudes can reinforce constructively or
cancel destructively, depending on phase.\status{E} This is the principle
behind all useful quantum algorithms: One prepares the register
in a uniform intermediate torsion over all inputs (Hadamard);
applies a function that increases the phases of the correct answers
and decreases those of the wrong ones; lets interference drive the wrong
answers into unstable configurations and the correct one into a
stable one; and readout drives the system into the nearest stable
pole, which after the interference corresponds with overwhelming frequency to the
sought answer. Without interference: Monte Carlo. With
interference: Grover ($\sqrt N$ instead of $N$ steps), Shor (factorization
in polynomial time).

In the FFGFT, the phase is not an abstract complex number, but the
rotational position of the intermediate torsion configuration. Constructive
interference: Two torsion phases lie equal and reinforce each other into a
geometrically more stable configuration. Destructive interference:
Two opposite phases cancel each other out; the resulting
configuration has no stable geometric anchor and is not
driven on readout. A quantum algorithm is, in this language, a sequence of torsion rotations that align the phase configuration
of the register so that the sought answer ends in a stable
configuration.

\begin{laierspur}
Water waves that meet reinforce each other when crest meets crest,
and cancel each other out when crest meets trough. A
quantum computer controls probabilities like waves: It lets the
correct answers grow and the wrong ones disappear. This only works
because quantum states have a phase that classical bits lack.
\end{laierspur}

\section{QND Measurement: Readout Without Destroying}

The readout interaction drives the intermediate configuration into
a pole --- for invasive measurement methods. There is a class of
methods that do not perform this step: the
quantum non-demolition measurement (QND). It reads out an observable $\hat A$
without changing the eigenstate; the formal condition is
$[\hat A,\hat H_{\mathrm{measurement}}] = 0$.\status{E} Practically: The same
state can be measured repeatedly and yields the same
result each time.

The most important example is Faraday or Kerr rotation: A linearly
polarized laser beam passes through the quantum dot without absorption and rotates its polarization plane by $\theta_F\propto S_z$. No photon is
absorbed, no energy transfer into the spin mode; the spin state is
identical before and after the measurement.

\textbf{QND as empirical evidence for determinism.} The argument is
direct: The Faraday rotation measures $S_z$ without changing the state;
repeated measurements yield the same result every time; therefore
the state had this value before every measurement, definitely, not as an abstract
superposition that only becomes reality through measurement. In the FFGFT, the Faraday rotation is state-preserving because a stable
torsion configuration cannot be driven out of its state by a dispersive, non-resonant coupling. The photon sees the
torsion --- it rotates its polarization accordingly --- but transfers no energy, because it is not resonant.\status{S}

\begin{center}
\small\renewcommand{\arraystretch}{1.3}
\begin{tabular}{p{2.4cm}p{3cm}p{2.8cm}}
\toprule
\textbf{Method} & \textbf{Interaction} & \textbf{State After} \\
\midrule
Resonant fluorescence & Photon absorbed, excites transition & changed \\
Spin-to-charge conversion & Electron tunnels out of the resonator & destroyed \\
Faraday/Kerr & Photon passes dispersively & unchanged: QND \\
\bottomrule
\end{tabular}
\end{center}

The standard reading, that the state only exists through measurement, is difficult to reconcile with QND measurements. In the FFGFT, the
torsion configuration is a geometric reality: stable, definite,
pre-existent. Measurement reads it off; it does not create it.

\textbf{Limits of QND.} QND is possible for observables that commute with the
measurement operator, typically $S_z$. An intermediate configuration
--- a state that is not an $S_z$ eigenstate --- is driven into an eigenstate by an
$S_z$ QND measurement, deterministically depending on the initial
state; the intermediate torsion is then gone.

\section{The Potential Landscape}

The complete picture of qubit states is a potential landscape
with two stable minima and a saddle in between. The minimum
$|\uparrow\rangle$ is a stable configuration, deep enough that
weak couplings cannot lift the state out of it --- QND possible. The
minimum $|\downarrow\rangle$ is mirror-symmetrically equally deep. The slope
is the intermediate torsion: The phase determines how far the state is
from which minimum; not a true minimum, metastable on
the coherence time $T_2$. Every point in between is a geometrically
defined state, continuously parameterized, discrete in the
endpoints.

The readout strength decides: A state in the stable minimum remains
under weak coupling (QND) and can be lifted out under strong coupling
(invasive). A state on the slope rolls into the nearest minimum even under weak
coupling --- which one, the exact phase determines,
epistemically unknown, deterministic; under strong coupling the same,
only faster.

The crucial insight: There is no principled boundary between
\enquote{QND possible} and \enquote{QND impossible}, only more stable and
less stable states in a continuous landscape. The jump into the final position is not
random: A state near $|\uparrow\rangle$ lands in $|\uparrow\rangle$, one near
$|\downarrow\rangle$ in $|\downarrow\rangle$, one on the saddle
depending on the fine phase. The apparent randomness of the
standard QM is the epistemic ignorance of this phase. The picture holds
for all $d$ levels of a qudit: $d$ minima with $d-1$ saddle points.

\section{Two Levels of the $\xi$ Hierarchy: Spin Qubit and Photonic Resonator}

The previous sections developed the qubit on the spin qubit ---
an electron in a nanometer cage. There is a qualitatively different
realization class: the photonic resonator. Both enforce discrete
states through geometry, but sit on different levels
of the $\xi$ hierarchy, with consequences for temperature, stability, and
the nature of information.

\textbf{Spin qubit: minimal resonator at level $N\approx8$.} A
single electron, one energy level, two orientations. The separation
of the spin states in a typical field of 1\,T is
$\Delta E_Z = g^*\mu_BB\approx0{,}1$--$1$\,meV; the thermal energy
at room temperature is $k_BT = 26$\,meV. Thus $\Delta E_Z\ll k_BT$:
Thermal noise constantly throws the system out of the minimum. Deep
temperatures below 1\,K are not optional, but mandatory --- not
to preserve superposition, but to remain in one of the minima at all.

\textbf{Photonic resonator: collective geometry at level
$N\approx9$--$10$.} A waveguide of lithium niobate, a
Mach-Zehnder interferometer, a ring resonator --- resonators of
$10^{18}$ to $10^{23}$ atoms in lattice order. The resonator is not
a single particle, but the macroscopic geometry of the entire
crystal; the photon belongs to the collective mode. The discreteness comes
from the boundary conditions of the macroscopic geometry. The energy of an
optical photon is $1$--$2$\,eV, thus $E/k_BT\approx40$: Thermal
noise cannot create optical photons. Photonics works
at room temperature not because it is built more robustly, but because
it operates on a higher $\xi$ level: The geometric robustness
is given by the single crystal with $10^{23}$ lattice sites, the thermal by the eV scale --- four orders of magnitude apart.

\textbf{Not limited to two states.} In a Mach-Zehnder
interferometer, the photon runs through two paths with continuously
adjustable phase; the interference result $\cos^2(\phi/2)$ is a
continuous transfer function. A photonic resonator can operate both modes:
digital with $\phi\in\{0,\pi\}$ and transmission
1 or 0, or analog with continuous phase --- infinitely many
intermediate states, all stable at room temperature. In the spin qubit,
the analogous information lies in a thermally unstable intermediate state,
$T_2$ collapses in nanoseconds at room temperature; in the
photonic resonator, it lies in the phase of the photon, and that is
thermally stable.

\begin{center}
\small\renewcommand{\arraystretch}{1.3}
\begin{tabular}{p{2.6cm}p{2.8cm}p{3cm}}
\toprule
 & \textbf{Spin Qubit} & \textbf{Photonic Resonator} \\
\midrule
$\xi$ level & $N\approx8$ (nm) & $N\approx9$--10 ($\mu$m--mm) \\
Resonator & 1 electron in cage & $10^{18}$--$10^{23}$ atoms collective \\
Energy scale & $0{,}1$--1\,meV & 1--2\,eV \\
$\Delta E/k_BT$(300\,K) & $0{,}004$--$0{,}04$ & $\approx40$ \\
Cooling needed? & yes, $T<1$\,K & no \\
Analog states & fragile ($T_2\sim$ns) & stable \\
Discreteness from & confinement (1 particle) & mode geometry \\
Typical application & digital qubit & analog signal processing, photonic computing \\
\bottomrule
\end{tabular}
\end{center}

Spin qubit and photonic resonator are, in the FFGFT, not
fundamentally different systems, but resonators at different
levels of the same scaling ladder. Each level carries the energy by the
factor $1/\xi\approx7500$ higher than the previous one; from meV (level 8) to
eV (level 9--10) the factor is $10^3$--$10^4$, consistent with the
scaling.\status{S} This also explains why biological systems
use photonic mechanisms (photosynthesis, bird navigation) and no
spin qubits for quantum effects: Nature works on the
thermally robust level.

\section{The Analog Resolution Limit}

The photonic resonator is thermally robust, but only as long as one
does not resolve too finely. Anyone who wants to distinguish analog intermediate values
divides the energy space into $N$ levels with $\Delta E_{\mathrm{level}} =
E_{\mathrm{photon}}/N$, and the condition $\Delta E_{\mathrm{level}}\gg
k_BT$ yields $N\ll E_{\mathrm{photon}}/k_BT$. At room temperature and
optical photons: $N_{\max}\approx40$, corresponding to
$\log_2 40\approx5{,}3$\,bit per photon. The photonic system does not escape
the thermal problem --- it shifts it to a finer
scale. More resolution requires cooling or an even higher
energy scale.

This is the oldest engineering wisdom of communications engineering,
geometrically grounded: High-frequency technology has been climbing the same
hierarchy for a hundred years. At each frequency level, the signal-to-noise
ratio improves because $E/k_BT$ grows with each level.

\begin{center}
\footnotesize\renewcommand{\arraystretch}{1.25}
\begin{tabular}{p{2.2cm}p{1.3cm}p{1.6cm}p{1.2cm}p{2.9cm}}
\toprule
\textbf{Technology} & \textbf{Frequency} & \textbf{$E/k_BT$} & \textbf{$N$-Level} & \textbf{Regime} \\
\midrule
AM radio & 1\,MHz & $1{,}5\times10^{-7}$ & $\approx5{,}5$ & thermally dominated \\
FM & 100\,MHz & $1{,}5\times10^{-5}$ & $\approx6{,}4$ & noise-limited \\
Microwave & 10\,GHz & $1{,}5\times10^{-3}$ & $\approx7{,}3$ & manageable \\
5G/mmWave & 30\,GHz & $5\times10^{-3}$ & $\approx7{,}5$ & transition \\
THz & 1\,THz & $0{,}15$ & $\approx7{,}8$ & $E\sim k_BT$: difficult \\
Near-infrared & 200\,THz & 31 & $\approx8{,}8$ & thermally robust \\
Fiber optic & 300\,THz & 46 & $\approx9{,}0$ & $N_{\max}\approx46$ \\
UV & 1000\,THz & 154 & $\approx9{,}4$ & $N_{\max}\approx150$ \\
\bottomrule
\end{tabular}
\end{center}

From $N = 6$ (AM) to $N = 9$ (optics), three levels bridge the factor
$(1/\xi)^3\approx4\times10^{11}$; in fact, the fiber optic frequency is
$10^8$-fold above AM radio. The fiber optic is, in this picture, not a
new technology, but the continuation of the same trend to the
next level. And the resolution limit applies at every level: Anyone who wants to resolve more than 40 analog levels at optical wavelengths must cool --- as the radio engineer in 1930 shielded his receiver. The
universal limit reads
\begin{equation}
  N_{\max}(T,N_\xi)\approx\frac{E_0\,(1/\xi)^{N_\xi-N_0}}{k_BT}.
\end{equation}

\textbf{Above optics.} Optics is not the end of the ladder,
but the end of the usable zone for guided photonic transmission.
In the near UV (3--10\,eV, $N\approx9{,}3$) photoionization begins: The
resonator is damaged by its own signal. At soft X-rays
(100\,eV--1\,keV, $N\approx10$) the wavelength approaches the
lattice spacing, Bragg scattering at individual lattice planes, the
waveguiding collapses. At hard X-rays, the photon
flies through the lattice gaps, matter becomes partially transparent, there is no
refractive index and no resonator. At gamma energies, the
photon interacts only with atomic nuclei --- another $\xi$ level. Above
1\,MeV ($N\approx11{,}3$) pair production sets in: the natural
upper limit. In the FFGFT there are no cannonballs and no waves as
separate concepts, only torsion configurations at various scales;
waveguiding is a scale resonance of two torsion patterns, the
crystal at level 7--8 and the photon at level 9, and it breaks down
when the scales become incompatible.

\begin{keypoint}[What This Chapter Holds]
The qubit is a mode in the cylindrical phase space; gates are
torsion rotations; registers grow exponentially through topological
connection; interference is phase superposition of torsions; QND measurement
evidences the pre-existence of stable poles; the potential landscape makes the
measurement outcome deterministic with epistemically unknown phase. Spin qubit
and photonics are resonators at different $\xi$ levels, and the
thermal resolution limit applies equally at every level.
\end{keypoint}

\quelle{Doc.~175 (Qubit state spaces, March 2026), Doc.~034 (cylindrical formalism), Doc.~173 (resonator conditions), Doc.~174 (torsion, Faraday readout), Doc.~230}